\section{Related Work}
\label{sec:related}

Our work sits at the intersection of benchmark contamination detection, contamination-resistant evaluation design, fine-tuning dynamics, and the emerging crisis of trust in LLM benchmarks. We organize prior work into five themes and identify the gap that motivates our study.

\subsection{Benchmark Contamination Detection}

The problem of detecting whether evaluation data has leaked into training corpora has received substantial attention. \textbf{Probability-based methods} pioneered by \citet{shi2024detecting} introduced Min-K\% Prob, which flags contamination by examining the average log-probability of the $k$\% least probable tokens in a sequence, under the hypothesis that unseen examples contain more low-probability outlier tokens. This was subsequently improved by \citet{zhang2025minkpp}, who provided theoretical grounding by framing detection as identification of local maxima in the modeled distribution. \citet{xie2024recall} proposed ReCaLL, a membership inference attack leveraging the insight that conditioning on non-member prefixes induces larger log-likelihood drops for member data. More recently, \citet{choi2025kernel} introduced the Kernel Divergence Score (KDS), which measures changes in kernel similarity matrices before and after fine-tuning on benchmark data, achieving near-perfect correlation with contamination levels.

\textbf{Prompt-based detection} approaches include \citet{golchin2024timetravel}'s ``guided instruction'' method, where models are prompted with dataset names and partial instances to test completion ability, achieving 92--100\% accuracy. \citet{deng2024investigating} introduced Testset Slot Guessing, revealing that GPT-4 can guess 57\% of missing MMLU options. \citet{yang2024rethinking} demonstrated that simple n-gram overlap methods miss rephrased contamination, proposing an LLM-based decontaminator that found 8--18\% HumanEval overlap in common pre-training corpora.

\textbf{N-gram and perplexity-based methods} remain widely used despite known limitations. \citet{zhou2024benchmarking} developed a pipeline combining perplexity and n-gram accuracy to analyze 31 LLMs, revealing substantial test set misuse in mathematical reasoning and proposing the Benchmark Transparency Card. \citet{simulating2025leakage} compared detection methods under realistic continual pre-training, finding n-gram methods achieve the highest F1 scores. In the emerging area of post-training detection, \citet{selfcritique2025rl} proposed Self-Critique, the first method targeting RL post-training contamination, achieving up to 30\% AUC improvement by probing for policy collapse.

A critical finding across this literature is that \textbf{no detection method operates across the full training pipeline}. Each targets a single stage---pre-training \citep{shi2024detecting,zhang2025minkpp}, supervised fine-tuning \citep{fsd2024finetuning}, or RL post-training \citep{selfcritique2025rl}---without tracking how contamination signals transform across stages. Our work fills this gap by measuring detection method efficacy at each pipeline stage for the same underlying contamination.

\subsection{Contamination-Resistant Benchmarks}

The complementary approach to detecting contamination is preventing it through benchmark design. \citet{white2025livebench} introduced LiveBench, which draws from recent math competitions, arXiv papers, and news articles with monthly updates, establishing the gold standard for temporal contamination resistance. \citet{jain2024livecodebench} applied similar principles to code evaluation, continuously collecting problems from competitive programming platforms. Domain-specific variants include LiveMedBench \citep{livemedbench2026} for clinical medicine, which revealed that 84\% of models degrade on post-cutoff cases.

An alternative strategy generates test instances algorithmically. \citet{beyondbench2025} introduced BeyondBench with over $10^{15}$ unique instances across 44 algorithmic tasks with deterministic verification, evaluating 101 models. \citet{dycodeval2025} uses multi-agent systems to generate semantically equivalent code problem variations, revealing significant performance degradation suggestive of contamination.

\citet{zhang2024gsm1k} created GSM1K, a human-authored mirror of GSM8K, finding up to 8\% accuracy drops and systematic overfitting in Mistral and Phi model families---a landmark result published at the NeurIPS 2024 Datasets and Benchmarks Track. \citet{wang2024mmlupro} developed MMLU-Pro with 10-option multiple choice and harder questions, producing 16--33\% accuracy drops relative to MMLU. \citet{antileakbench2025} constructed AntiLeak-Bench using Wikidata/Wikipedia updates to guarantee contamination-free evaluation. In the software engineering domain, \citet{lessleak2025} audited 83 benchmarks, finding generally low leakage (4.8\% average for Python) but extreme cases such as QuixBugs at 100\%.

While these benchmarks address contamination \emph{for future evaluation}, they cannot retroactively assess whether \emph{existing fine-tuned models} carry contamination inherited from their base models---the question we address.

\subsection{Contamination Propagation and Forgetting}

The most directly related line of work studies how contamination behaves under continued training. \citet{bordt2025forgetting} systematically studied forgetting along three scaling dimensions (parameters up to 1.6B, repetitions up to 144$\times$, training tokens up to 40B), finding that contamination can be ``forgotten'' when training data scales beyond 5$\times$ Chinchilla---a regime characteristic of modern LLMs. However, their study examines only the pre-training phase and does not study what happens when fine-tuning is applied on top of potentially contaminated knowledge.

\citet{zhang2026leakagetrap} is the closest existing work to ours, studying contamination propagation through LoRA adaptation in LLM-based recommendation systems. They identified a ``Triple Effect'' where domain-relevant leakage creates spurious performance gains while domain-irrelevant leakage degrades accuracy. Their study is limited to the recommendation domain and a single LoRA configuration on a single model family; our study differs by covering two independent model families, four distinct downstream task types, three contamination rates, and strict exact-match scoring on seven general-purpose benchmarks.

Two recent papers study the interaction between reinforcement learning and memorization. \citet{wu2025reasoning} demonstrated that Qwen2.5's apparent mathematical reasoning gains from random or incorrect rewards are artifacts of pre-training contamination on MATH-500 and AIME benchmarks, proposing the uncontaminated RandomCalculation alternative. \citet{yan2026spurious} provided mechanistic evidence via Path Patching and Logit Lens, identifying an ``Anchor-Adapter circuit'' (layers 18--20 anchoring memorized solutions, layers 21+ adapting representations) that enables memorization shortcuts during RLVR training. These findings are compelling but limited to the Qwen model family and mathematical reasoning tasks.

\textbf{Our work differs} from all of the above in three key respects. First, we study contamination propagation through pre-training followed by clean supervised fine-tuning on four \emph{unrelated} downstream tasks, whereas \citet{bordt2025forgetting} examine forgetting only during continued pre-training and \citet{zhang2026leakagetrap} examine only in-domain LoRA fine-tuning. Second, we evaluate on seven general-purpose benchmarks (GSM8K, MMLU, HumanEval, MATH, MBPP, ARC-Challenge, TruthfulQA-MC) across two independent model families, rather than a single benchmark family or domain. Third, we measure whether four existing contamination detection methods remain effective at identifying contamination that has been transformed by downstream supervised fine-tuning, which no prior study has directly quantified. We do not evaluate RLHF or DPO post-training in this paper; that scope boundary is flagged explicitly as a limitation in Section~\ref{sec:discussion}.

\subsection{Decontamination Methods}

Several approaches attempt to mitigate contamination effects. \citet{zhu2024itd} proposed Inference-Time Decontamination (ITD), which detects and rewrites leaked samples at evaluation time, reducing inflated accuracy by 22.9\% on GSM8K and 19.0\% on MMLU. \citet{chai2026deconIEP} extended this idea with DeconIEP, applying bounded perturbations in embedding space guided by less-contaminated reference models.

However, \citet{sun2025emperor} provided a sobering assessment: testing 20 mitigation strategies across 10 LLMs and 5 benchmarks, they found that \textbf{no existing strategy effectively balances fidelity and contamination resistance}. Semantic-preserving strategies offered no improvement over doing nothing, while semantic-altering strategies sacrificed benchmark validity. \citet{xu2025dcr} proposed a more nuanced approach with the DCR framework, quantifying contamination at four granular levels (semantic, informational, data, label) and producing adjusted accuracy scores within 4\% of uncontaminated baselines. For code-specific applications, \citet{codecleaner2025} developed automated refactoring operators achieving 75\% overlap reduction.

Our work reveals a further complication for decontamination: even if base model contamination is mitigated, fine-tuning may reactivate or transform residual contamination signals, necessitating pipeline-aware decontamination strategies.

\subsection{Fine-Tuning Dynamics and Knowledge Retention}

Understanding how fine-tuning affects pre-trained knowledge is essential background for our work. \citet{wu2025forgetting} showed that training on LLM-generated data (which exhibits lower perplexity) reduces non-target task degradation, proposing Selective Token Masking to achieve similar benefits with ground-truth data. \citet{upweight2025easy} found that upweighting samples with low pre-trained model loss preserves knowledge during SFT. \citet{empirical2024forgetting} documented that catastrophic forgetting intensifies with scale (1B--7B), and notably that LoRA does \emph{not} mitigate forgetting in continual learning---contradicting common assumptions.

At the intersection of forgetting and safety, \citet{leveraging2024forgetting} showed catastrophic forgetting can be leveraged positively for safety alignment, while \citet{factual2024pretraining} established a power-law relationship between training steps and factual knowledge forgetting during pre-training.

These works study forgetting of \emph{general knowledge}, but none specifically examine whether \emph{memorized benchmark answers} are forgotten, retained, or amplified during supervised fine-tuning---the precise question our paper answers for the SFT stage.

\subsection{Benchmark Trust Crisis}

Our work is further motivated by the growing crisis of trust in LLM evaluation. The Llama~4 controversy, where Meta submitted a customized model variant to LMArena that ranked 2nd while the public release ranked 32nd---later confirmed by Yann LeCun as ``fudged'' results---exemplifies the institutional pressures driving benchmark manipulation. \citet{li2025preference} identified ``preference leakage'' in LLM-as-a-judge paradigms, where judges systematically favor outputs from related models, compounding contamination concerns. \citet{eriksson2025trust} conducted a meta-review of approximately 100 studies on benchmark shortcomings, cataloging issues from dataset biases to contamination to signal-noise confusion, concluding that the benchmark ecosystem faces a systemic validity crisis. Multiple surveys \citep{deng2024unveiling,xu2024bdcsurvey,li2025dcsurvey,comprehensive2025detection,static2dynamic2025} have documented the growing scale of the contamination problem, with \citet{assumptions2024detection} finding that 3 of 8 common detection assumptions fail across training phases.

Our paper contributes to resolving this crisis by demonstrating a previously uncharacterized failure mode: contamination in open-weight base models silently propagates to thousands of downstream fine-tuned variants, creating a systemic contamination problem that no amount of careful benchmark construction can fully address without awareness of the propagation mechanism.
